% Compilar: cd docs && latexmk -pdf barras.tex && latexmk -c
% Figuras: scripts_barras/fig_*.py (astro-clean); cada cifra del texto sale de lo que imprimen.
\documentclass[11pt,a4paper]{article}
\usepackage[utf8]{inputenc}
\usepackage[T1]{fontenc}
\usepackage[spanish,es-noshorthands]{babel}
\usepackage{amsmath}
\usepackage[varg]{txfonts}
\usepackage{graphicx}
\usepackage{booktabs}
\usepackage{tabularx}
\usepackage[margin=2.5cm]{geometry}
\usepackage[hidelinks]{hyperref}
\decimalpoint

\newcommand{\as}{\ensuremath{^{\prime\prime}}}
\newcommand{\aeps}{a_\varepsilon}
\newcommand{\code}[1]{\texttt{#1}}

\title{Medir el tamaño de una barra en DESI DR1\\[2pt]
\large Isofotas, Fourier y descomposición 2D frente al método de Alonso et al. (2024)}
\author{Frank Bautista}
\date{15 de noviembre de 2026}

\begin{document}
\maketitle

\begin{abstract}
% Se explica cómo se mide el largo, el ancho y la orientación de una barra en las imágenes de DESI Legacy, sin
% aprendizaje automático. El largo y el ángulo salen del ajuste de isofotas (el radio donde la elipticidad es
% máxima); el ancho, de la descomposición de Fourier de la luz, que separa la barra (modo $m = 2$) de la luz
% redonda del bulbo y el disco ($m = 0$). Se compara con el método de Alonso et al. (2024), que fija un contorno
% de brillo a partir de cuatro puntos marcados a mano y entrena redes para imitarlos. En las 448 barras que midió
% Valeria, el largo por isofotas coincide con el suyo (mediana del cociente $1.01$, el 88\,\% a menos del 20\,\%),
% también en las 291 que no se usaron para ajustar el método. El ancho depende de la definición: el de las
% isofotas y el de Valeria incluyen el bulbo y coinciden ($0.97$); el de Fourier no lo incluye y es la mitad. Se
% documentan los errores encontrados por el camino, cómo se corrigieron y qué queda pendiente: barras de pocos
% segundos de arco, la deproyección y las máscaras.
Se explica cómo se mide el largo, el ancho y la orientación de una barra en las imágenes $r$ de DESI Legacy, sin
aprendizaje automático. El largo y el ángulo salen del ajuste de isofotas (el primer pico de la elipticidad,
buscado más allá de $\max(2\as, 1\,{\rm kpc})$); el ancho, de la descomposición de Fourier de la luz, que separa
la barra del bulbo y el disco. El disco de las isofotas externas da la deproyección, validada con S$^4$G. El error
del largo es el tramo donde la elipticidad no se distingue del pico (9\,\% en mediana); el ruido de la imagen
aporta menos del 2\,\%. En las 448 barras de Alonso et al. (2024), el largo coincide con el suyo (mediana del
cociente $1.01$, el 88\,\% a menos del 20\,\%) y las diferencias tienen el tamaño del error. La
sección~\ref{sec:metodo} resume el método; el resto explica por qué, con los errores encontrados y lo pendiente.
\end{abstract}

\tableofcontents


\section{El método que se usa}\label{sec:metodo}

Esta sección resume lo que funciona; el resto del documento explica por qué. Después de las pruebas de las
secciones~\ref{sec:isofotas}--\ref{sec:resultados}, el método base para correr muestras grandes es:
\begin{enumerate}
\item \textbf{Banda $r$.} Es la banda de todo el documento y es más profunda; $z$ empata (sección~\ref{sec:bandas}).
\item \textbf{Isofotas y primer pico de $\varepsilon$} con $\varepsilon \ge 0.25$ y una caída posterior de 0.1, buscado
desde
\begin{equation}\label{eq:amin}
a_{\min} = \max\!\left(2\as,\ \frac{1\ \mathrm{kpc}}{s(z)}\right),
\end{equation}
con $s(z)$ la escala en kpc por segundo de arco. Da $a = \aeps$ y PA $=$ PA$(\aeps)$.
\item \textbf{Ancho de Fourier}, $W$ a $0.5\,\aeps$ y $0.75\,\aeps$ (barra combinada, ecuación~\ref{eq:comb}).
\item \textbf{Disco y deproyección.} $\varepsilon_d$ es la mediana de $\varepsilon$ y PA$_d$ el promedio (en ángulo doble)
de las isofotas con $23.5 < \mu_r < 25$\,mag\,arcsec$^{-2}$ y $a > 1.5\,\aeps$ (si son menos de tres, las tres
últimas). Se guardan $a$, $W$ y $L$ proyectados y deproyectados, $\varepsilon_d$, PA$_d$, $\Delta\phi$ y $a_{\min}$.
\item \textbf{Bandera $r/z$.} Se mide también en $z$ y se marca la galaxia si $|a_r/a_z - 1| > 0.15$
(sección~\ref{sec:bandas}). $z$ no reemplaza a $r$: solo avisa.
\end{enumerate}
Lo implementan \code{measure\_image} y \code{measure\_galaxy} en \code{src/measure.py}; la bandera $r/z$ todavía no
está en el código.

\subsection{El error de la medida}\label{sec:error}

Se reporta \textbf{un solo error por medida}: el tramo de la \emph{meseta} del pico. Son las isofotas contiguas
al pico con $\varepsilon \ge \varepsilon_{\rm pico} - \sigma_\varepsilon$, donde $\sigma_\varepsilon$ es el error de
$\varepsilon$ en el pico; su extremo interior y exterior dan $[a_{\rm lo}, a_{\rm hi}]$, con un piso de medio paso
entre isofotas. Es asimétrico porque el perfil lo es. En las 448:
\begin{itemize}
\item $\sigma_a/a$ tiene mediana 0.09 (16--84\,\%: 0.05--0.13).
\item \textbf{Está calibrado:} la diferencia con Valeria dividida por este error tiene dispersión 1.09 (lo
esperado es 1), y Valeria cae dentro de $[a_{\rm lo}, a_{\rm hi}]$ en el 60\,\% de las galaxias (68\,\% esperado).
No está inflado; si algo, se queda un poco corto.
\item \textbf{Marca las dudosas:} con $\sigma_a/a > 0.3$ hay 8 galaxias y el 62\,\% falla frente a Valeria
(11\,\% en el resto). Se usa como bandera ``revisar''.
\end{itemize}

\paragraph{El ruido de la imagen es menor.} Se sumó a la imagen $r$ ruido gaussiano con la varianza de DESI
(\code{ivar}) 20 veces y se volvió a medir, en 100 de las 448. Como la imagen ya tiene ruido, la dispersión se
divide por $\sqrt{2}$. El error por ruido es 1.5\,\% de $a$ (16--84\,\%: 0.6--2.5\,\%) y 1.6\,\% de $W$; en 16 de
las 100 el pico no se mueve nunca y solo en 4 salta más de un 20\,\%. Queda contenido en la meseta y basta
mencionarlo.

\paragraph{Lo que se descartó.} (1) Monte Carlo sobre el perfil (ruido independiente en cada $\varepsilon$): crea
picos falsos antes de la barra, porque el error de isofotas vecinas está correlacionado. (2) El paso entre isofotas
como atajo del error por ruido: es tres veces mayor que éste. (3) La dispersión entre definiciones del largo (caída
de $\varepsilon$, giro del PA o de la fase de $m=2$, caída de $A_2$): no es un error sino una cota superior; en la
galaxia ejemplo todas caen en $18.7$--$19.9\as$, por fuera del pico ($15.4\as$) y de Valeria ($14.05\as$).

\subsection{Qué necesita una lista de galaxias}\label{sec:lista}

Para medir basta la posición. \textbf{RA y DEC} son obligatorios: con ellos se baja el recorte de Legacy y todo
lo demás (cielo, máscara, centro, isofotas, Fourier, disco) sale de la imagen. \textbf{$z$ es opcional}: sin él,
$a$, $W$ y PA salen en segundos de arco y $a_{\min} = 2\as$, pero no hay $L$ en kpc ni mínimo de 1\,kpc. El
\code{SHAPE\_R} de Tractor solo fija el lado del recorte (8 radios efectivos, entre 128 y 1024 píxeles; en la
muestra de 27\,283, la mediana es 173 píxeles). Sin él hay que elegir el lado de otra forma (pendiente).

\subsection{Por qué un mínimo en kpc}

En galaxias cercanas $2\as$ es muy poco: el primer pico de $\varepsilon$ es una estructura nuclear resuelta (anillo,
lente o barra nuclear, de pocos cientos de pc) y la barra queda corta. Se probó con 30 barras de S$^4$G
(Herrera-Endoqui et al. 2015) medidas en DESI con este método, con escalas de $\sim 0.14$\,kpc por segundo de arco:
\begin{center}
\begin{tabular}{lccc}
\toprule
$R_{\min}$ & S$^4$G a menos del 20\,\% (de 30) & Mediana / S$^4$G & 448: a menos del 20\,\% \\
\midrule
solo $2\as$ & 7 & 0.62 & 88.4\,\% \\
0.75\,kpc & 12 & 0.81 & 88.4\,\% \\
1\,kpc & 14 & 0.93 & 88.4\,\% \\
1.25\,kpc & 13 & 0.95 & 88.6\,\% \\
2\,kpc & 11 & 1.20 & 88.6\,\% \\
\bottomrule
\end{tabular}
\end{center}
Con 1\,kpc se duplican los aciertos en S$^4$G y ninguna de las 448 cambia: están a $z = 0.012$--$0.093$, donde
$2\as$ ya es más de 1\,kpc (el mínimo en kpc solo actúa a $z < 0.025$). La barra más corta de Valeria mide 1.96\,kpc,
así que el límite no quita barras reales; por encima de 1.5--2\,kpc sí empieza a saltárselas. Las que siguen
fallando en S$^4$G son sobre todo barras de galaxias tardías (Sd, Sm), donde la propia medida visual de S$^4$G y la
de su perfil de elipticidad difieren hasta en un factor 2.

A $z = 0.1$, en cambio, el límite es $2\as = 3.8$\,kpc: no se pueden medir barras más cortas. Por eso se guarda
$a_{\min}$ de cada galaxia, para tener en cuenta ese sesgo al comparar muestras.

\subsection{La deproyección, validada con S$^4$G}

El \code{SHAPE} de Tractor no sirve para el disco: en las barradas sigue a la barra (difiere de ella en
$4.7^\circ$ en mediana). Las isofotas externas sí: en las 448 el factor de deproyección
$a_{\rm dep}/a$ tiene mediana 1.09. En las mismas 30 de S$^4$G, el factor que se obtiene con el disco medido en DESI
dividido por el que da la orientación de S$^4$G tiene mediana 0.98 (16--84\,\%: 0.90--1.04) y el 93\,\% queda a
menos del 20\,\%.


\section{Por qué medir barras}\label{sec:contexto}

Este contexto sigue las clases de Virginia Cuomo (Formación y Evolución de Galaxias, ULS, 2025, módulo III).

\paragraph{Qué son.} Las barras se conocen desde la clasificación de Hubble (1926); de Vaucouleurs (1959) separó
las familias SA (sin barra), SAB (barra débil) y SB (barra fuerte). Suelen ir acompañadas de anillos, brazos que
nacen de sus puntas, barras internas y bulbos en forma de caja o maní. Una barra no es un objeto estático: rota
como un cuerpo rígido con una \textbf{velocidad de patrón} $\Omega_p$ y lleva gas hacia el centro, donde forma
estructuras nuevas. Por eso domina la evolución secular del disco: intercambia momento angular, se frena y crece.

\paragraph{Cómo se forman.} Dos escenarios (Cuomo, clase 2, p.\,52--55; clase 3, p.\,8--9):
\begin{enumerate}
\item \textbf{En aislamiento}, por inestabilidad del disco cuando es dinámicamente frío.
\item \textbf{Inducidas por una interacción} con otra galaxia.
\end{enumerate}
La pregunta abierta es si la barra ``recuerda'' cómo se formó. La forma de saberlo es el \textbf{cociente de
rotación}
\begin{equation}\label{eq:R}
\mathcal{R} = \frac{R_{\rm cr}}{R_{\rm bar}} = \frac{V_{\rm circ}}{\Omega_p\,R_{\rm bar}},
\end{equation}
donde $R_{\rm cr}$ es el radio de corrotación. Las barras \emph{rápidas} tienen $1 \le \mathcal{R} \le 1.4$ y las
\emph{lentas} $\mathcal{R} > 1.4$; las simulaciones indican que las barras inducidas por interacciones son lentas
(Martínez-Valpuesta et al. 2017). $\Omega_p$ y $V_{\rm circ}$ requieren cinemática (método de Tremaine--Weinberg),
pero $R_{\rm bar}$ es \textbf{fotometría}: es lo que mide este trabajo, y es un ingrediente obligado de
$\mathcal{R}$.

\paragraph{Relación con los pares.} Es la misma pregunta que motiva a Alonso et al. (2024) y a la tesis: si las
interacciones en pares cercanos cambian las barras. Con el largo, el ancho y la elipticidad de miles de barras en
DESI se puede comparar pares contra control. Dos precauciones de las clases:
\begin{itemize}
\item \textbf{Tendencias conocidas} (Cuomo et al. 2020): las galaxias más masivas o brillantes tienen barras más
largas, y las barras más largas son más fuertes. Pares y control deben compararse a igual masa.
\item \textbf{Tamaño relativo}: en el cúmulo de Virgo las barras del campo son más largas, pero las galaxias del
campo también son más grandes, y el largo \emph{normalizado por el tamaño de la galaxia} es el mismo (clase 3,
p.\,28--29). Hay que comparar también $R_{\rm bar}$ normalizado, como hace el paper con el tamaño del disco.
\end{itemize}

\paragraph{Fuerza de la barra.} La fuerza mide cuánto contribuye la barra al potencial de la galaxia. Una
estimación fotométrica directa es la elipticidad de la barra, $\varepsilon_{\rm bar}$ (Aguerri et al. 2009):
barras más alargadas son más fuertes. Este trabajo ya la mide (\code{eps\_bar}).

\paragraph{No hay un borde nítido.} El semilargo de una barra es la extensión de sus órbitas estelares, pero la
barra no tiene un borde definido en la imagen (Hilmi et al. 2020). Por eso hay varios métodos fotométricos
(sección~\ref{sec:isofotas}) y dan valores que difieren entre sí; ninguno es ``el verdadero''. Lo importante es
usar un criterio publicado y aplicarlo igual a toda la muestra.


\section{Qué se quiere medir}\label{sec:problema}

Una barra es una estructura alargada y recta en el centro de un disco. Se describe con tres números: su
\textbf{semilargo} $a$ (del centro a la punta), su \textbf{semiancho} $b$ y su \textbf{ángulo de posición} PA.
El largo total es $L = 2a$ y el ancho $W = 2b$. Su forma se resume con la elipticidad
\begin{equation}
\varepsilon_{\rm bar} = 1 - \frac{b}{a}.
\end{equation}

\subsection{De segundos de arco a kpc}

En la imagen se mide un ángulo $\theta$. Su tamaño físico, \emph{proyectado sobre el plano del cielo}, es
\begin{equation}\label{eq:kpc}
L\,[\mathrm{kpc}] = D_A(z)\,\theta = \frac{D_M(z)}{1+z}\,\theta,
\qquad
D_M(z) = \frac{c}{H_0}\int_0^{z}\frac{dz'}{E(z')},\quad E(z) = \sqrt{\Omega_m(1+z)^3 + \Omega_\Lambda},
\end{equation}
con $\theta$ en radianes ($1\as = 4.848\times10^{-6}$\,rad) y la cosmología Planck18. Es la misma conversión
que la separación proyectada de un par de galaxias,
\begin{equation}\label{eq:rp}
r_p = \frac{D_M(\bar z)\,\theta}{1+\bar z},
\end{equation}
donde $\theta$ es la separación angular de los dos miembros. Para la galaxia ejemplo
($z = 0.039$) la escala es $0.807$\,kpc por segundo de arco.

\subsection{Proyectado y deproyectado}

Lo que se pierde al proyectar es distinto en los dos casos. En un par, la separación a lo largo de la línea de
visión es desconocida y no se recupera: $r_p$ es una cota inferior de la separación real. En una barra, lo que se
pierde es la \emph{inclinación del disco}, $i$, y eso sí se puede corregir, porque la barra vive en el plano del
disco. Si el disco tiene elipticidad aparente $\varepsilon_d$ y ángulo PA$_d$, entonces $\cos i = 1 - \varepsilon_d$
(disco delgado), y un radio a lo largo de la barra, que forma un ángulo
$\Delta\phi = \mathrm{PA}_{\rm bar} - \mathrm{PA}_d$ con el eje mayor del disco, se deproyecta como
\begin{equation}\label{eq:deproj}
a_{\rm dep} = a\,\sqrt{\cos^2\Delta\phi + \frac{\sin^2\Delta\phi}{\cos^2 i}} .
\end{equation}
% El código deproyecta la elipse completa de la barra (escribe la elipse como forma cuadrática, estira la
% dirección del eje menor del disco en $1/\cos i$ y saca los nuevos semiejes de los autovalores), pero la
% orientación del disco todavía no es confiable (sección~\ref{sec:pendientes}). Por eso aquí, igual que en el
% paper, \textbf{todo se reporta proyectado}.
El método guarda \textbf{las dos medidas}: la proyectada, para comparar con Valeria (el paper no deproyecta), y la
deproyectada con \eqref{eq:deproj}, para comparar con trabajos que sí lo hacen. El ancho se deproyecta con la misma
fórmula y el ángulo $90^\circ - \Delta\phi$, porque es perpendicular a la barra. Cómo se mide el disco y su
validación están en la sección~\ref{sec:metodo}.


\section{Definiciones: la elipse y la barra}\label{app:definiciones}

La figura~\ref{fig:definiciones} reúne las cantidades que aparecen en el documento.

\paragraph{Una isofota (figura~\ref{fig:definiciones}a).}
\begin{itemize}
\item $(x_0, y_0)$: el centro de la elipse, en píxeles. El código lo fija en el centro de la galaxia.
\item $a$: el semieje mayor. Es el eje horizontal de todos los perfiles $\varepsilon(a)$, PA$(a)$ y $\mu(a)$.
\item $b$: el semieje menor, $b = a(1-\varepsilon)$.
\item $\varepsilon = 1 - b/a$: la elipticidad. Vale 0 para un círculo y tiende a 1 para una elipse muy alargada. A
veces se usa la razón de ejes $q = b/a = 1 - \varepsilon$.
\item $\theta$: el ángulo del eje mayor en el sistema de la imagen, medido desde $+x$ en sentido antihorario. Es el
que entrega \code{photutils} y el de los perfiles de la figura~\ref{fig:isofotas}c.
\item PA: el ángulo de posición astronómico, medido desde el Norte hacia el Este. En el cielo el Este queda a la
izquierda, así que también va en sentido antihorario. Es el que reporta el catálogo (\code{pa\_bar}):
$\mathrm{PA} = \theta - 90^\circ$, módulo $180^\circ$. Un eje no tiene sentido, por lo que $0^\circ$ y
$180^\circ$ son la misma orientación.
\item $r(\varphi)$: la distancia del centro al borde de la elipse en una dirección que forma un ángulo $\varphi$
con el eje mayor,
\begin{equation}
r(\varphi) = \frac{a\,b}{\sqrt{(b\cos\varphi)^2 + (a\sin\varphi)^2}}, \qquad r(0) = a, \quad r(90^\circ) = b .
\end{equation}
Solo sobre el eje mayor la distancia al borde es $a$. Dos elipses con distinto ángulo pueden tener $a$ muy distinto
y llegar casi a la misma distancia en una dirección dada (sección~\ref{sec:figuras}).
\end{itemize}

\paragraph{La barra (figura~\ref{fig:definiciones}b).}
\begin{itemize}
\item $\aeps$: el $a$ de la isofota en el primer pico de $\varepsilon$. Es el semilargo de la barra, $a_{\rm bar}$;
su elipticidad y su ángulo son $\varepsilon_{\rm bar}$ y PA$_{\rm bar}$.
\item $a_{\rm up}$: el $a$ de la primera isofota después del pico con $\varepsilon$ menor en 0.1 o con el PA girado
más de $10^\circ$. Es una cota superior del largo: esa isofota ya pertenece al disco o al anillo.
\item $b_{\rm iso} = \aeps(1-\varepsilon_{\rm bar})$: el semieje menor de la isofota de la barra. Incluye la luz
del bulbo (sección~\ref{sec:isofotas}).
\item $W$: el ancho \emph{total} de la barra, medido en la imagen solo-barra de Fourier (modos $m = 2, 4, 6$) como el
FWHM de los cortes perpendiculares en $0.5\,a_{\rm bar}$ y $0.75\,a_{\rm bar}$. La caja se dibuja con $W/2$ hacia
cada lado.
\item $L = 2\,a_{\rm bar}$: el largo total. En kpc, $L_{\rm kpc} = L \times$ escala$(z)$ y $W_{\rm kpc} = W
\times$ escala$(z)$, con la escala en kpc por segundo de arco para Planck18.
\end{itemize}
Las salidas del código son $a_{\rm bar}$ (semilargo) y $W$ (ancho total), ambas proyectadas en el cielo.

\begin{figure}[t]
\centering
\includegraphics[width=\textwidth]{figures_barras/definiciones.pdf}
\caption{Definiciones. (a) Una isofota: centro, semiejes $a$ y $b$, ángulo $\theta$ desde $+x$ (sistema de la
imagen), PA desde el Norte hacia el Este (astronómico) y distancia al borde $r(\varphi)$ en la dirección $\varphi$.
(b) La barra con su eje sobre $x$: isofotas grises que se alargan hacia afuera, la de $\aeps$ (azul), con
$a_{\rm bar}$ y $b_{\rm iso}$; la barra medida (morado), con su ancho total $W$; la isofota de $a_{\rm up}$
(naranja a trazos), ya girada y casi redonda; y el largo total $L$. Esquema, sin datos.}
\label{fig:definiciones}
\end{figure}


\section{El método de Alonso et al. (2024)}\label{sec:paper}

Alonso et al. (2024) midieron las barras de 300 galaxias (150 en pares y 150 de control) sobre imágenes JPG de
SDSS de $500\times500$ píxeles a $0.30\as$/px:
\begin{enumerate}
\item Identifican \emph{a ojo} el ángulo de la barra y giran la imagen para dejarla sobre el eje $x$.
\item Atenúan el canal azul, suavizan con un núcleo gaussiano, pasan a gris y normalizan entre 0 y 1.
\item Marcan \emph{a mano} cuatro puntos: A y B en las puntas de la barra sobre el perfil en $x$, C y D en sus
bordes sobre el perfil en $y$, los dos perfiles por el centro.
\item Promedian el brillo de los cuatro puntos, $I_c = \tfrac14(I_A + I_B + I_C + I_D)$, y toman el contorno
$I(x,y) = I_c$ como la barra: su extensión en $x$ es $2a$ y en $y$ es $2b$.
\item Entrenan una ResNet50 para reproducir el ángulo (error $\sim 2^\circ$) y una red densa que, a partir de los
perfiles en $x$ e $y$, predice el nivel $I_c$. Con ellas miden el resto de la muestra.
\end{enumerate}
El paper \textbf{no deproyecta}: limita la muestra a galaxias de cociente de ejes $b/a > 0.4$ (casi de frente),
solo gira la imagen para dejar la barra sobre $x$ y reporta el cociente de ejes de la barra como
\emph{proyectado}. Las redes \emph{imitan} el criterio de quien marcó: no hay una definición física del fin de la barra. El paper no
compara con métodos estándar ni discute el efecto del bulbo.

La figura~\ref{fig:paper} reproduce el procedimiento en la galaxia ejemplo con los extremos que midió Valeria
($a = 14.05\as$, $b = 4.83\as$). El contorno resultante tiene semiejes $13.8\as$ y $4.85\as$, pero no es una
barra: es una S que sigue el arranque de los brazos. Además, C y D están sobre el perfil \emph{que pasa por el
centro}, donde domina el bulbo: el ancho mide hasta dónde el bulbo y el disco son más brillantes que $I_c$, no el
grosor de la barra.

\begin{figure}[t]
\centering
\includegraphics[width=\textwidth]{figures_barras/paper.pdf}
\caption{Método de Alonso et al. (2024) en la galaxia ejemplo (TARGETID 39627835815496281, $z = 0.039$), sobre
la imagen FITS de DESI suavizada. (a) Imagen girada con la barra sobre $x$, los cuatro puntos y el contorno al
nivel promedio (naranja). (b) Perfiles por el centro en $x$ (arriba) y en $y$ (abajo), con el nivel promedio.}
\label{fig:paper}
\end{figure}


\section{Ajuste de isofotas}\label{sec:isofotas}

\subsection{Qué se ajusta}

Una isofota es la curva de brillo constante. El método de Jedrzejewski (1987), implementado en \code{photutils},
la aproxima por una elipse de centro $(x_0, y_0)$, semieje mayor $a$, elipticidad $\varepsilon$ y ángulo PA (definiciones en la sección~\ref{app:definiciones}).
Sobre la elipse se muestrea el brillo en función de un ángulo $\varphi$ que recorre la elipse, medido desde su
eje mayor, y se ajusta
\begin{equation}\label{eq:jedr}
I(\varphi) = I_0 + A_1\sin \varphi + B_1\cos \varphi + A_2\sin 2\varphi + B_2\cos 2\varphi .
\end{equation}
No confundir $\varphi$ con $\theta$: $\theta$ es la orientación de la elipse (un número por isofota), y
$\varphi$ la posición sobre ella (va de $0^\circ$ a $360^\circ$ al dar la vuelta).
Si la elipse sigue bien una isofota, el brillo no cambia a lo largo de ella y los cuatro coeficientes son cero.
Cada uno indica cómo corregir la geometría: $A_1, B_1$ mueven el centro, $A_2$ corrige el ángulo y $B_2$ la
elipticidad, divididos por el gradiente radial $\partial I/\partial a$. Se itera hasta que el coeficiente mayor
es pequeño frente al ruido. Repitiendo para $a = 1, 1.05, 1.05^2,\dots$ píxeles se obtienen los perfiles
$\varepsilon(a)$, PA$(a)$ e $I(a)$ (figura~\ref{fig:isofotas}). La sección~\ref{app:isofotas} detalla el procedimiento paso a paso.

Las elecciones del código: el centro se fija; el brillo sobre cada elipse es la \emph{mediana} (robusta a regiones
HII, polvo y estrellas) con recorte a $3\sigma$; y se acepta un gradiente con error relativo de hasta 2 (el valor
por defecto, 0.5, detiene el ajuste en la barra, donde el perfil es casi plano).

\subsection{Leer la barra en las curvas}

Una barra tiene una firma clara (figura~\ref{fig:isofotas}b,c). En el centro $\varepsilon$ es baja porque el
bulbo es casi redondo; a lo largo de la barra $\varepsilon$ sube con PA casi constante; en su extremo
$\varepsilon$ cae y PA gira hacia el disco o los brazos. Las clases de Cuomo (clase 2, p.\,21--22) resumen los
tres métodos fotométricos estándar para el radio de la barra:
\begin{enumerate}
\item \textbf{Máximo de la elipticidad}: el radio del máximo local de $\varepsilon$ (Wozniak \& Pierce 1991). Es
el que se usa aquí para el largo.
\item \textbf{Cambio del ángulo}: el radio donde el PA cambia $5^\circ$ respecto al del máximo de
$\varepsilon$. Aquí se usa $10^\circ$ (Erwin 2005) y solo para la bandera \code{flag\_end} (``la barra termina''),
no para el largo; adoptar $5^\circ$ no cambiaría ningún largo.
\item \textbf{Análisis de Fourier}: los modos pares tienen un pico donde hay barra, y la extensión del pico da el
radio (sección~\ref{sec:fourier}). Aquí se usa para el ancho.
\end{enumerate}
El ejercicio de la clase 2 (p.\,27) es exactamente el método 1 sobre una galaxia: $R_{\rm bar} = 11\as$ del
máximo de $\varepsilon$. Los criterios que se aplican:
\begin{itemize}
\item \textbf{Hay barra} si $\varepsilon$ alcanza un máximo $\varepsilon_{\max}\ge 0.25$ y después cae al
menos $\Delta\varepsilon = 0.1$ (Marinova \& Jogee 2007). Se toma el \emph{primer} pico que cumple esto, no el
máximo global, porque un disco inclinado o los brazos externos pueden ser más alargados que la barra.
\item \textbf{Largo} (Erwin 2005): la cota inferior es el radio del máximo, $\aeps$; la superior, $a_{\rm up}$, es
el primer radio posterior donde $\varepsilon < \varepsilon_{\max} - 0.1$ o el PA se aparta más de $10^\circ$.
\end{itemize}
El semiancho de la barra por isofotas es el semieje menor de esa elipse,
\begin{equation}\label{eq:biso}
b_{\rm iso} = \aeps\,(1-\varepsilon_{\max}) ,
\end{equation}
y el ángulo, el de la isofota en $\aeps$. Para el ejemplo: $\aeps = 15.4\as$, $\varepsilon_{\max} = 0.67$,
$a_{\rm up} = 18.7\as$, $b_{\rm iso} = 5.1\as$; Valeria midió $14.05\as$. No se busca la barra por debajo de
$a_{\min} = 2\as$ (unas 1.5 veces el seeing en banda $r$); este umbral es provisorio.

\begin{figure}[t]
\centering
\includegraphics[width=\textwidth]{figures_barras/isofotas.pdf}
\caption{Ajuste de isofotas en la galaxia ejemplo. (a) Una de cada cinco elipses ajustadas (gris) y la de la
barra (azul). (b) Elipticidad, (c) ángulo de posición y (d) brillo superficial frente al semieje mayor. Línea
azul continua: $\aeps$; punteada: $a_{\rm up}$; negra a trazos: semilargo de Valeria. Gris: $a < a_{\min}$.}
\label{fig:isofotas}
\end{figure}

\subsection{El defecto del ancho}

El largo se mide en la punta de la barra, lejos del bulbo. El ancho no: el semieje menor de la elipse
\eqref{eq:biso} tiene su vértice sobre el corte perpendicular \emph{que pasa por el centro}. Allí la isofota
envuelve la luz de la barra, del bulbo y del disco, que son más redondos. Esa mezcla diluye la elipticidad
($\varepsilon_{\max}$ subestima la de la barra) y $b_{\rm iso}$ sale inflado (Gadotti 2008). Es el mismo problema
que el punto C y D del paper. Además, la isofota real en $\aeps$ no es una elipse sino una S, barra más arranque
de los brazos (figura~\ref{fig:contornos}a), y la elipse que la resume se alarga y ensancha.


\subsection{El ajuste, paso a paso}\label{app:isofotas}

Esta subsección detalla lo que hace \code{src/isophotes.py} con \code{photutils} (Jedrzejewski 1987). La
figura~\ref{fig:iteracion} sigue el ajuste de una isofota real de la galaxia ejemplo, en $a = 5\as$.

\begin{figure}[t]
\centering
\includegraphics[width=\textwidth]{figures_barras/iteracion.pdf}
\caption{Ajuste de una isofota de la galaxia ejemplo en $a = 5\as$, con el centro fijo. (a) Elipse de partida
(naranja a trazos, $\varepsilon = 0.30$, $\theta = 140^\circ$) y ajustada (azul). (b) Brillo a lo largo de la
elipse de partida: mediana por sector (gris), armónicos 1 y 2 de la ecuación~\eqref{eq:armonicos} (naranja) e
$I_0$ (negro). (c) Lo mismo para la elipse ajustada, después de 10 iteraciones. $\theta$ se mide en el sistema de la
imagen, desde $+x$; en (b) y (c) el eje horizontal es $\varphi$, la posición sobre la elipse medida desde su
eje mayor (lo que devuelve \code{photutils}), no $\theta$.}
\label{fig:iteracion}
\end{figure}

\paragraph{Paso 0. La idea.} Una isofota es la curva donde el brillo es constante, como una curva de nivel en un
mapa topográfico. El método supone que cada isofota se parece a una elipse y busca, para cada tamaño, la elipse
que mejor la sigue. Una elipse tiene cinco parámetros: el centro $(x_0, y_0)$, el semieje mayor $a$, la
elipticidad $\varepsilon = 1 - b/a$ y el ángulo $\theta$ de su eje mayor (en el sistema de la imagen, desde $+x$).
El semieje $a$ no se ajusta: el método recorre los radios uno a uno. El código además fija el centro, así que en
cada radio solo se ajustan $\varepsilon$ y $\theta$.

\paragraph{Paso 1. La elipse de prueba.} Para un $a$ dado se parte de una geometría de prueba. El punto de la
elipse en la dirección $\varphi$, el ángulo polar medido desde su eje mayor, está a la distancia
\begin{equation}
\begin{gathered}
r(\varphi) = \frac{b}{\sqrt{(1-\varepsilon)^2\cos^2\varphi + \sin^2\varphi}}, \qquad b = a(1-\varepsilon), \\
x = x_0 + r\cos(\varphi + \theta), \qquad y = y_0 + r\sin(\varphi + \theta).
\end{gathered}
\end{equation}
Así lo hace \code{photutils} (\code{EllipseGeometry.radius}). Jedrzejewski usa en cambio el ángulo excéntrico
$E$, el de las órbitas de Kepler: $x' = a\cos E$, $y' = b\sin E$ en el sistema de la elipse, con
$\tan\varphi = (1-\varepsilon)\tan E$. Los dos coinciden sobre los ejes y difieren entre ellos; en un círculo
($\varepsilon = 0$) son el mismo.
En la figura~\ref{fig:iteracion} la elipse de partida (naranja a trazos) tiene $\varepsilon = 0.30$ y
$\theta = 140^\circ$.

\paragraph{Paso 2. El brillo a lo largo de la elipse.} Se recorre la elipse y se mide el brillo $I(\varphi)$.
Con \code{integrmode="median"} el anillo elíptico alrededor de la elipse se divide en sectores, y en cada sector
se toma la mediana de los píxeles cuyo centro cae dentro. Si un sector tiene 6 píxeles o menos, lo que pasa en
radios chicos, \code{photutils} interpola de forma bilineal entre los cuatro píxeles vecinos. Después se
descartan los valores a más de $3\sigma$ de la media, tres veces, para eliminar estrellas, regiones HII y
píxeles enmascarados. El resultado son los puntos grises de la figura~\ref{fig:iteracion}b.

\paragraph{Paso 3. ¿Es constante?} Si la elipse siguiera exactamente una isofota, $I(\varphi)$ sería una línea
horizontal. Si oscila, la elipse está mal puesta: cruza zonas más brillantes y más débiles. En la
figura~\ref{fig:iteracion}b el brillo oscila dos veces por vuelta, entre 0.24 y 0.34\,nmgy/px.

\paragraph{Paso 4. Los armónicos.} Las oscilaciones se describen con una serie de Fourier corta en $\varphi$,
ajustada por mínimos cuadrados:
\begin{equation}\label{eq:armonicos}
I(\varphi) = I_0 + A_1\sin\varphi + B_1\cos\varphi + A_2\sin 2\varphi + B_2\cos 2\varphi .
\end{equation}
$I_0$ es el brillo medio de la isofota; $A_1, B_1$ describen una oscilación por vuelta y $A_2, B_2$ dos. Es la
ecuación~\eqref{eq:jedr}. Jedrzejewski la escribe con $E$ y \code{photutils} con $\varphi$; para lo que sigue
la diferencia no importa. En la partida,
$A_2 = 0.022$ y $B_2 = -0.034$\,nmgy/px, grandes frente a la dispersión de los puntos.

\paragraph{Paso 5. Por qué cada armónico indica un error distinto.} Este es el centro del método. Sea
$\delta r(\varphi)$ la distancia radial entre la elipse y la isofota verdadera en la dirección $\varphi$. Como el
brillo cambia con el radio según el gradiente $I' = \partial I/\partial a$ (negativo: el brillo baja hacia
afuera), a primer orden en $\delta r$
\begin{equation}\label{eq:taylor}
I(\varphi) \approx I_0 + I'\,\delta r(\varphi) .
\end{equation}
El brillo medido es el gradiente por cuánto se equivocó la elipse de lugar, y cada error deja una huella
distinta en $\varphi$:
\begin{center}
\begin{tabular}{lll}
\toprule
Error de la elipse & $\delta r(\varphi)$ & Armónico \\
\midrule
centro corrido sobre el eje mayor & $\propto\cos\varphi$ & $B_1$ \\
centro corrido sobre el eje menor & $\propto\sin\varphi$ & $A_1$ \\
ángulo girado & $\propto\sin 2\varphi$ & $A_2$ \\
elipticidad equivocada & $\propto\cos 2\varphi$ & $B_2$ \\
\bottomrule
\end{tabular}
\end{center}
Por ejemplo, si la elipse es más redonda que la isofota, se sale de ella en las dos puntas del eje mayor
($\varphi = 0^\circ$ y $180^\circ$) y queda dentro en los costados: dos veces por vuelta, con forma de coseno,
es decir $B_2$. Si está girada, el desfase es de $45^\circ$ y aparece $A_2$.

\paragraph{Paso 6. La corrección.} Despejando $\delta r$ de la ecuación~\eqref{eq:taylor} sale cuánto mover cada
parámetro (las fórmulas de \code{photutils.isophote.fitter}):
\begin{equation}\label{eq:correcciones}
\delta\varepsilon = -\frac{2B_2(1-\varepsilon)}{a\,I'}, \qquad
\delta\theta = \frac{2A_2(1-\varepsilon)}{a\,I'\,\left[(1-\varepsilon)^2 - 1\right]}, \qquad
\delta x_{\rm mayor} = -\frac{B_1}{I'}, \qquad
\delta x_{\rm menor} = -\frac{A_1(1-\varepsilon)}{I'} .
\end{equation}
En cada iteración se corrige \emph{solo} el parámetro cuyo armónico es mayor en valor absoluto, y se vuelve al
paso~2 con la elipse nueva. Con el centro fijo, $A_1$ y $B_1$ se miden pero no se usan. El gradiente se mide con
una segunda elipse un 5\,\% más grande,
\begin{equation}
I' \approx \frac{I_0(1.05\,a) - I_0(a)}{0.05\,a} ,
\end{equation}
y su error sale de la dispersión del brillo sobre las dos elipses.

\paragraph{Paso 7. Cuándo se detiene.} El ajuste converge (código 0) cuando el armónico mayor queda por debajo
del 5\,\% de la dispersión de los residuos respecto de la ecuación~\eqref{eq:armonicos}, después de al menos 10
iteraciones. Si no, sigue hasta 50 y devuelve la iteración con el armónico más chico (código 2). Otros códigos:
1 si se descartó más del 30\,\% de los puntos de la elipse, 3 si la elipse no tiene puntos válidos o el ajuste de armónicos falla,
y $-1$ si el gradiente no es válido. Solo los códigos 0 y 2 se usan en los perfiles; el código agrega el 6
(\code{STUCK}) para los códigos 2 que terminaron con la elipticidad de partida sin tocar. En el ejemplo, tras 10
iteraciones queda $\varepsilon = 0.21$, $\theta = 158^\circ$ (figura~\ref{fig:iteracion}c): $A_2$ y $B_2$
bajaron a 0.001 y la oscilación pasó de 0.10 a 0.03\,nmgy/px. Queda $A_1 = 0.015$ porque el centro está fijo; el
resto son brazos y polvo, que la elipse no puede seguir.

\paragraph{Paso 8. El punto débil: el gradiente.} Todas las correcciones~\eqref{eq:correcciones} se dividen por
$I'$. Si el brillo casi no cambia con el radio, $I' \approx 0$ y las correcciones se disparan. Eso pasa justo en
la barra, cuyo perfil es casi plano (la meseta de la figura~\ref{fig:cortes}b). Por eso \code{photutils} rechaza
la isofota si el error relativo del gradiente supera \code{maxgerr}; el valor por defecto, 0.5, detenía el ajuste
en la barra, y el código usa 2. Aun así, una elipse de partida lejana no converge en la barra: con la misma
partida de la figura~\ref{fig:iteracion} en $a = 7$--$12\as$, el ajuste se aleja y termina con código $-1$.

\paragraph{Paso 9. De radio en radio.} Por lo anterior, cada radio parte de la última elipse aceptada, que ya
está cerca de la verdadera y solo necesita correcciones chicas. El código empieza en $a = 15$ píxeles
($3.9\as$) con la elipticidad y el ángulo de los momentos de segundo orden de la luz, y avanza hacia afuera con
$a_{k+1} = 1.05\,a_k$ hasta el borde de la galaxia. Después recorre hacia adentro, hasta 1 píxel, partiendo de la
primera isofota que convergió. Cada radio se ajusta por separado: la rutina \code{fit\_image} de \code{photutils}
congela la geometría de todos los radios mayores cuando dos isofotas seguidas tienen gradiente nulo o positivo,
y una barra o un anillo bastan para que eso pase. El resultado es una tabla con $a$, $I_0$, $\varepsilon$,
$\theta$, sus errores y el código de parada para cada radio: los perfiles de la figura~\ref{fig:isofotas}.




\section{Perfiles de brillo sobre la barra}\label{sec:figuras}

En las clases de Cuomo y en muchos trabajos aparecen figuras parecidas a la figura~\ref{fig:paper}b: el brillo
a lo largo de la barra y perpendicular a ella. Pero se usan para otra cosa que en el paper. Esta sección explica
cómo se construyen, qué se lee en ellas y en qué se diferencia el uso del paper. Ninguna de las curvas es un
ajuste a los puntos: son medidas directas sobre la imagen. El único ajuste es el del disco exponencial.

\subsection{Cortes por los ejes de la barra}

La construcción (figura~\ref{fig:cortes}):
\begin{enumerate}
\item Se gira la imagen FITS sin cielo para dejar la barra sobre el eje $x$, con el PA de $\aeps$.
\item \textbf{Corte por el eje mayor}: se promedian 5 filas de píxeles alrededor del centro (franja azul), para
bajar el ruido sin mezclar luz de fuera de la barra. \textbf{Corte por el eje menor}: lo mismo con 5 columnas
(franja morada).
\item Se promedian los dos lados del centro, de modo que cada corte queda en función del radio, y el brillo se
pasa a magnitudes: $\mu_r = 22.5 - 2.5\log_{10}(I/s^2)$, con $I$ en nanomaggies por píxel y $s = 0.262\as$/px.
\item Como referencia se agregan $\mu(a)$ de las isofotas (el brillo medio sobre cada elipse) y un disco
exponencial, $\mu(r) = \mu_0 + 1.086\,r/h$, ajustado por mínimos cuadrados a $\mu(a)$ entre $25\as$ y $42\as$,
fuera de la barra y del anillo. En el ejemplo, $h = 17\as$.
\item El \textbf{contraste barra/interbarra} es la diferencia de los dos cortes, $\mu_{\rm menor} - \mu_{\rm
mayor}$: cuántas magnitudes más brillante es la barra que el disco a la misma distancia del centro.
\end{enumerate}

Lo que se lee:
\begin{itemize}
\item \textbf{El tipo de barra} (Elmegreen \& Elmegreen 1985). Si el corte mayor tiene una meseta, la barra es
\emph{plana}, típica de barras fuertes en galaxias tempranas. Si cae de forma exponencial, como el disco, es una
barra \emph{exponencial}, más común en galaxias tardías. El ejemplo tiene una meseta entre $\sim 7\as$ y $13\as$:
barra plana.
\item \textbf{La fuerza de la barra}: el contraste dentro de ella. En el ejemplo vale $\sim 0.9$\,mag entre
$7\as$ y $13\as$, es decir, a lo largo de la barra hay unas 2.4 veces más luz que de través.
\item \textbf{Dónde termina}: donde el contraste cae a cero, cuando el corte mayor baja al nivel del corte menor
(Ohta et al. 1990; el cociente de Aguerri et al. 2000 es la misma idea con modos de Fourier). Con un umbral de
0.25\,mag, elegido aquí y no estándar, el fin queda en $17.6\as$. Es otra cota superior, como $a_{\rm up}$: marca
el fin del exceso de luz, que está más afuera que el pico de $\varepsilon$ ($\aeps = 15.4\as$; Valeria,
$14.05\as$).
\item \textbf{Cuánto domina sobre el disco}: el exceso del corte mayor sobre la recta del disco. Más allá de
$\sim 25\as$ los dos cortes y $\mu(a)$ siguen el disco; las oscilaciones del contraste allí son los brazos y el
anillo cruzando las franjas.
\end{itemize}
Este análisis describe la barra, pero no da un largo más preciso que el de las isofotas: el fin del contraste
depende del umbral y de los brazos. Por eso el código no lo usa para medir.

\begin{figure}[t]
\centering
\includegraphics[width=\textwidth]{figures_barras/cortes.pdf}
\caption{Cortes de brillo por los ejes de la barra en la galaxia ejemplo. (a) Imagen con la barra sobre $x$ y
las franjas de 5 píxeles de cada corte. (b) Brillo superficial a lo largo del eje mayor (azul) y del menor
(morado a trazos), $\mu(a)$ de las isofotas (naranja) y el disco exponencial ajustado entre $25\as$ y $42\as$
(verde). (c) Contraste barra/interbarra. Líneas verticales: $\aeps$ (azul), fin del contraste a 0.25\,mag (gris
punteada) y semilargo de Valeria (negra a trazos).}
\label{fig:cortes}
\end{figure}

\subsection{El mismo corte en el paper (figura~\ref{fig:paper}b)}

El paper usa los mismos dos cortes, pero no para mirar la forma de la barra: los usa para fijar \emph{un nivel de
brillo}. Las diferencias de construcción son que el corte es de un solo píxel de ancho, que la imagen es el gris
del JPG normalizado entre 0 y 1 (aquí se reproduce con el FITS suavizado por una gaussiana de $\sigma = 1.5$
píxeles, en nanomaggies y escala logarítmica) y que no se pasa a magnitudes.
\begin{itemize}
\item \textbf{Puntos A, B, C, D}: en el paper se marcan \emph{a mano} sobre los cortes, A y B en las puntas de la
barra y C y D en sus bordes. En la figura~\ref{fig:paper} no se marcan a ojo: se ponen en los extremos que midió
Valeria ($x = \mp 14.05\as$, $y = \mp 4.83\as$) y se lee el brillo del corte en esos píxeles.
\item \textbf{Línea roja a trazos}: el nivel $I_c$. No es un ajuste ni se elige a ojo: es el promedio aritmético
de los cuatro brillos, $I_c = \tfrac14(I_A + I_B + I_C + I_D)$. Por construcción los cuatro puntos quedan cerca
de ella.
\item \textbf{Contorno de la figura~\ref{fig:paper}a}: la curva de la imagen suavizada donde el brillo vale
exactamente $I_c$. Su extensión en $x$ e $y$ es $2a$ y $2b$.
\end{itemize}
La parte subjetiva es la posición de los cuatro puntos; la red neuronal aprende a predecir $I_c$ a partir de los
dos cortes. La diferencia con la sección anterior es de fondo: allí el fin de la barra sale de comparar los dos
cortes entre sí y con el disco; aquí sale de un único nivel de brillo, que en el eje menor cae sobre el bulbo.

\subsection{Perfiles de isofotas (figura~\ref{fig:isofotas}b--d)}

Cada punto es \emph{una} elipse ajustada (sección~\ref{sec:isofotas}). El eje horizontal es su semieje mayor
$a$, que crece en pasos de 5\,\% ($a_{k+1} = 1.05\,a_k$); por eso los puntos se juntan cerca del centro. Solo se
dibujan las elipses que convergieron.
\begin{itemize}
\item \textbf{(b) $\varepsilon$}: la elipticidad de cada elipse, $\varepsilon = 1 - b/a$.
\item \textbf{(c) PA}: el ángulo de cada elipse en el sistema de la imagen (desde $+x$, antihorario), entre
$0^\circ$ y $180^\circ$. Un eje no tiene sentido, así que $0^\circ$ y $180^\circ$ son la misma orientación: el
salto de $\sim 5^\circ$ a $\sim 175^\circ$ cerca de $a = 33\as$ es un giro pequeño, no uno de $170^\circ$. El
catálogo reporta el PA astronómico, desde el Norte hacia el Este: $\mathrm{PA} = \mathrm{PA}_{\rm imagen} -
90^\circ$, módulo $180^\circ$.
\item \textbf{(d) $\mu_r$}: el brillo superficial, a partir de la mediana $I$ del brillo sobre la elipse, en
nanomaggies por píxel: $\mu_r = 22.5 - 2.5\log_{10}(I/s^2)$, con $s = 0.262\as$/px la escala del píxel. El eje
está invertido (más brillante arriba). La meseta entre $\sim 6\as$ y $15\as$ es la barra; después del fin de
la barra el brillo vuelve a caer.
\item \textbf{Líneas verticales}: azul continua en $\aeps$, el primer pico de $\varepsilon$; azul punteada en
$a_{\rm up}$, la primera elipse después del pico con $\varepsilon$ menor en 0.1 o PA girado más de $10^\circ$;
negra a trazos en el semilargo de Valeria. La banda gris es $a < a_{\min}$, donde no se busca la barra.
\end{itemize}
En la galaxia ejemplo, $a_{\rm up}$ cumple las dos condiciones a la vez: entre $18.0\as$ y $18.7\as$ la
elipticidad baja de 0.65 a 0.21 y el PA gira $35^\circ$. El anillo corta la barra de golpe, y cualquier umbral de
PA entre $5^\circ$ y $35^\circ$ da el mismo $a_{\rm up}$.


\section{La descomposición de Fourier}\label{sec:fourier}

\subsection{La idea}

Si se recorre un círculo de radio $r$ alrededor del centro, en un disco o un bulbo redondos el brillo es
constante; si el círculo cruza una barra, el brillo sube dos veces por vuelta, una a cada lado
(figura~\ref{fig:fourier}b). En cada radio se escribe el brillo como serie de Fourier en el ángulo $\varphi$:
\begin{equation}\label{eq:fourier}
I(r,\varphi) = A_0(r) + \sum_{m=1}^{6}\left[a_m(r)\cos m\varphi + b_m(r)\sin m\varphi\right],
\end{equation}
con amplitud y fase de cada modo
\begin{equation}\label{eq:amfase}
I_m(r) = \sqrt{a_m^2 + b_m^2},\qquad \varphi_m(r) = \frac{1}{m}\arctan\frac{b_m}{a_m}.
\end{equation}
$A_0$ es la luz redonda (bulbo y disco). La barra, simétrica bajo un giro de $180^\circ$, vive en los modos pares:
$m = 2$ da su orientación y $m = 4, 6$ afinan su forma (puntas más rectangulares). $A_2/A_0 = I_2/A_0$ mide su
fuerza. Sin píxeles enmascarados, los coeficientes son la transformada de Fourier discreta de los $N = 360$
puntos del círculo,
\begin{equation}
a_m = \frac{2}{N}\sum_{n} I_n\cos m\varphi_n,\qquad b_m = \frac{2}{N}\sum_{n} I_n\sin m\varphi_n ;
\end{equation}
el código los obtiene por mínimos cuadrados sobre los puntos válidos, para que una estrella enmascarada no los
sesgue. No es una FFT de la imagen: es una serie en el ángulo, círculo por círculo.

\subsection{El largo y el ancho}

En la galaxia ejemplo, a $r = 8\as$ (sobre la barra) $A_2/A_0 = 0.45$ con fase $164^\circ$; a $r = 25\as$
(brazos) $A_2/A_0 = 0.51$ pero la fase es $74^\circ$: los brazos también tienen $m = 2$, incluso más que la
barra, pero su fase gira con el radio porque son espirales. De ahí dos criterios (Aguerri et al. 2000):
\begin{itemize}
\item \textbf{Hay barra} si el \emph{primer} pico de $A_2/A_0$ más allá de $a_{\min}$ supera 0.2.
\item \textbf{Largo por la fase}: la barra es recta, así que $\varphi_2$ es constante a lo largo de ella; termina
donde $|\varphi_2 - \varphi_2^{\rm pico}| > 10^\circ$. En el ejemplo, $a_{\varphi} = 11.8\as$
(figura~\ref{fig:fourier}e).
\end{itemize}
El criterio del cociente barra/interbarra de Aguerri et al. (2000),
$I_b/I_{ib} = (A_0 + I_2 + I_4 + I_6)/(A_0 - I_2 + I_4 - I_6)$, resultó inestable con anillos y brazos y no se usa.

Para el ancho se reconstruye la \emph{imagen solo-barra}, la ecuación \eqref{eq:fourier} sin $A_0$ y solo con
$m = 2, 4, 6$, dentro del radio de la barra (figura~\ref{fig:fourier}f). Como el bulbo es redondo, vive entero en
$A_0$ y desaparece por construcción. Se gira la imagen con la barra sobre $x$ y se mide el ancho a media altura
(FWHM) en cortes perpendiculares a $0.5$ y $0.75$ del largo, a cada lado,
\begin{equation}\label{eq:wf}
W_F = \frac14\sum_{k\in\{0.5,\,0.75\}}\ \sum_{s=\pm1} \mathrm{FWHM}\big[I_{\rm barra}(x = s\,k\,a,\ y)\big].
\end{equation}
No se corta en el centro porque allí los modos $m \ge 2$ valen cero (el ``corbatín'' de la figura). En el
ejemplo, $W_F = 5.6\as$, frente a $2b_{\rm iso} = 10.2\as$ de las isofotas.

\begin{figure}[t]
\centering
\includegraphics[width=\textwidth]{figures_barras/fourier.pdf}
\caption{Fourier en la galaxia ejemplo. (a) Círculos de medida. (b, c) Brillo de los píxeles a lo largo de cada
círculo (gris), $A_0$ (negro), $A_0$ más $m = 2$ (naranja) y la suma hasta $m = 6$. (d) $A_2/A_0$ y (e) fase
$\varphi_2$ frente al radio; naranja: fin de la barra por la fase; azul: $\aeps$; negra a trazos: Valeria.
(f) Imagen solo-barra ($m = 2, 4, 6$) girada, con la barra de Fourier.}
\label{fig:fourier}
\end{figure}


\section{Isofotas y Fourier: la misma herramienta, dos análisis}\label{app:comparacion}

Los dos métodos escriben el brillo como una serie de senos y cosenos en un ángulo, ecuaciones~\eqref{eq:armonicos}
y~\eqref{eq:fourier}. Pero no son el mismo análisis con más o menos términos: lo que cambia es sobre qué curva se
mide y qué se hace con los coeficientes.

\begin{center}
\small
\begin{tabularx}{\textwidth}{>{\raggedright\arraybackslash}p{3.1cm}XX}
\toprule
 & Isofotas (Jedrzejewski 1987) & Fourier (Aguerri et al. 2000) \\
\midrule
Pregunta & ¿qué forma tiene la curva de brillo constante? & ¿cuánto y cómo varía el brillo al dar la vuelta? \\
Curva de medida & elipse $(a, \varepsilon, \theta)$ que se deforma & círculo de radio $r$, fijo \\
Ángulo de la serie & $\varphi$ sobre la elipse, desde su eje mayor ($E$ en Jedrzejewski) & $\varphi$ sobre el círculo, desde $+x$ \\
Serie & $I_0 + \sum_{n=1}^{2}(A_n\sin n\varphi + B_n\cos n\varphi)$ & $A_0 + \sum_{m=1}^{6}(a_m\cos m\varphi + b_m\sin m\varphi)$ \\
Papel de los armónicos & errores de la elipse: $n = 1$ el centro, $n = 2$ $\varepsilon$ y $\theta$ & la estructura: $m = 0$ lo redondo, $m = 2$ la barra (y los brazos), $m = 4, 6$ su forma \\
Qué se hace con ellos & se llevan a cero moviendo la elipse & se guardan: son la medida \\
Cálculo & iterativo, 10 a 50 veces por radio; divide por el gradiente $I'$ & un ajuste lineal (mínimos cuadrados) por círculo \\
Puede fallar & sí: no converge si $I' \approx 0$ (meseta de la barra) o la partida está lejos & no: siempre da coeficientes \\
Muestreo radial & $a_{k+1} = 1.05\,a_k$ & cada píxel en $r$ \\
Máscaras & recorte $3\sigma$ y mediana por sector & se omiten los puntos enmascarados del ajuste \\
Salida & $\varepsilon(a)$, $\theta(a)$, $I_0(a)$ & $A_m(r)/A_0(r)$, $\varphi_m(r)$ \\
Firma de la barra & $\varepsilon$ con un máximo y $\theta$ constante & $A_2/A_0$ alto y $\varphi_2$ constante \\
Bulbo & queda dentro de la isofota: no se separa & queda en $m = 0$: se puede quitar \\
Brazos & la isofota pasa a ser del disco y $\varepsilon$ cae: no confunden el largo & también son $m = 2$: confunden el largo, no el ancho \\
Se usa para & largo $\aeps$ y ángulo & ancho $W$ (imagen solo barra, a $0.5$ y $0.75\,\aeps$) \\
\bottomrule
\end{tabularx}
\end{center}

Las isofotas preguntan qué forma tiene la curva de brillo constante: deforman la elipse hasta que el brillo
deja de variar, y la respuesta es la forma. Fourier pregunta cuánto varía el brillo al dar la vuelta a un círculo y
con qué patrón: deja la curva redonda, y la respuesta es la variación. Las dos ven la barra de maneras
complementarias. En las isofotas aparece como elipses alargadas ($\varepsilon$ alta, PA constante); en Fourier,
como un modo $m = 2$ grande con fase constante. Por eso cada una se usa para lo que hace mejor: las isofotas para el
largo y el ángulo; Fourier para el ancho, porque deja el bulbo, que es redondo, en $m = 0$ y lo separa de la barra.

Sobre la elipse ya ajustada, \code{photutils} calcula además los armónicos $n = 3$ y 4, que miden cuánto se aparta
la isofota de una elipse (forma de caja o de rombo). No cambian la geometría y el código no los usa.

Dos diferencias de notación con otros trabajos: el término constante a veces se escribe $A_0/2$ en vez de $A_0$
(es solo un factor 2), y cada autor nombra los coeficientes a su manera. En la ecuación~\eqref{eq:armonicos}
$A_n$ acompaña al seno; en otros textos acompaña al coseno. Lo que importa es que el coseno mide la onda alineada
con el ángulo cero y el seno, la desfasada.


\section{La barra combinada}\label{sec:combinada}

Fourier da un ancho sin bulbo, pero su largo por la fase se pasa en muchas galaxias (sección~\ref{sec:resultados}).
Las isofotas dan un largo que coincide con Valeria, pero un ancho inflado. La barra que se usa combina lo mejor de
cada uno:
\begin{equation}\label{eq:comb}
a = \aeps,\qquad \mathrm{PA} = \mathrm{PA}(\aeps),\qquad W = W_F\ \text{medido a } 0.5\,\aeps \text{ y } 0.75\,\aeps .
\end{equation}
El ángulo es el de la isofota en $\aeps$: promediarlo entre $\aeps/2$ y $\aeps$ (v4,
sección~\ref{sec:errores}) casi no lo cambia y la descomposición 2D favorece el de la isofota
(sección~\ref{sec:2d}). El ancho se evalúa directamente con la serie \eqref{eq:fourier} en los puntos de cada
corte, sin construir ni girar la imagen solo-barra. El ancho combinado de la galaxia ejemplo es $7.2\as$
(figura~\ref{fig:contornos}c). La figura muestra por qué hace
falta: la isofota real (a) es una S y la de Fourier (b) son dos lóbulos sin bulbo; ninguna es ``la silueta'' de la
barra. Los dos métodos entregan números ($a$, $W$, PA) y las elipses son solo una forma de dibujarlos.

\begin{figure}[t]
\centering
\includegraphics[width=\textwidth]{figures_barras/contornos.pdf}
\caption{Galaxia ejemplo con la barra sobre $x$. (a) Isofota real al nivel de $\aeps$ (azul) y su elipse (a
trazos). (b) Contorno a media altura de la imagen solo-barra. (c) Barras con los parámetros de cada método:
isofotas (azul), Fourier (naranja) y combinada (morado); círculos: extremos de Valeria.}
\label{fig:contornos}
\end{figure}


\section{Descomposición 2D}\label{sec:2d}

\subsection{El modelo}

GALFIT (Peng et al. 2002) e imfit (Erwin 2015; aquí \code{pyimfit}) ajustan a la imagen una suma de componentes
analíticas convolucionada con la PSF, minimizando
\begin{equation}\label{eq:chi2}
\chi^2 = \sum_{\text{píxeles}}\frac{\left[I_{\rm obs} - (\mathrm{PSF}\ast I_{\rm modelo})\right]^2}{\sigma^2},
\end{equation}
con $\sigma^2$ de la varianza inversa del recorte y la PSF del coadd de Legacy. Componentes: bulbo Sérsic,
\begin{equation}
I(r) = I_e\exp\left\{-b_n\left[\left(r/r_e\right)^{1/n}-1\right]\right\};
\end{equation}
disco exponencial, $I(r) = I_0\,e^{-r/h}$; anillo gaussiano; y barra Ferrers sobre una elipse generalizada,
\begin{equation}\label{eq:ferrers}
I(m) = I_0\left[1-\left(\frac{m}{a_{\rm bar}}\right)^2\right]^{n},\qquad
m = \left(|x|^{c_0+2} + \left|\frac{y}{1-\varepsilon}\right|^{c_0+2}\right)^{1/(c_0+2)},
\end{equation}
con $n = 2$; $c_0 > 0$ hace la barra más rectangular. Los valores iniciales salen de las isofotas.

\subsection{Qué aporta}

En la galaxia 39627884146461249, donde isofotas y Fourier no coincidían en el ángulo, el modelo sin anillo da
$\chi^2_\nu = 1.71$ y estira la barra hasta absorber el anillo; con anillo, $\chi^2_\nu = 1.20$
(figura~\ref{fig:2d}). La barra 2D tiene PA $= 133.8^\circ$: confirma el ángulo de la isofota en $\aeps$
($130.8^\circ$) y no el de Fourier ($120.5^\circ$), que mide su fase a $4\as$, donde todavía pesa el bulbo. En 10
galaxias de Valeria al azar, el PA 2D y el de las isofotas difieren $1.8^\circ$ en mediana.

El largo, en cambio, no es comparable: $a_{\rm bar}$ es donde el perfil \eqref{eq:ferrers} llega a cero, no donde
la barra deja de verse, y sale $1.47$ veces $\aeps$ en las 10 galaxias. En la imagen menos bulbo, disco y anillo
(figura~\ref{fig:2d}d) la barra visible termina cerca de Valeria ($10.9\as$). En 10 galaxias ciegas con barras de
$2$--$5\as$, el ajuste falla en 2 ($\chi^2_\nu = 562$ y $58$) y $c_0$ queda a menudo en sus límites. La
descomposición sirve como árbitro del ángulo en barras grandes, no como medida del largo.

\begin{figure}[h!]
\centering
\includegraphics[width=\textwidth]{figures_barras/descomposicion.pdf}
\caption{Descomposición de 39627884146461249. (a) Barras combinada (morado), de Fourier (naranja) y 2D (verde);
círculo: Valeria. (b) Modelo bulbo + disco + barra + anillo. (c) Residuo normalizado. (d) Imagen menos bulbo,
disco y anillo, con la barra 2D (verde) y la combinada (morado).}
\label{fig:2d}
\end{figure}

\subsection{Cómo leer la figura~\ref{fig:2d}}

La figura es una prueba de control, no parte del método. En esta galaxia las isofotas y Fourier daban ángulos
distintos para la barra, y hacía falta un tercer método, independiente, que hiciera de árbitro. La
descomposición no lee la imagen: construye un \emph{modelo} de la galaxia completa sumando piezas analíticas
(bulbo, disco, barra y anillo) y ajusta todos sus parámetros a la vez, hasta que el modelo convolucionado con la
PSF se parezca lo más posible a la imagen (ecuación~\ref{eq:chi2}).
\begin{itemize}
\item \textbf{(a)} La imagen con las tres barras: combinada (largo y ángulo por isofotas, ancho por Fourier),
solo Fourier, y la del modelo 2D.
\item \textbf{(b)} El modelo: la galaxia ``sintética'' que encontró el ajuste.
\item \textbf{(c)} El residuo, $(I_{\rm obs} - I_{\rm modelo})/\sigma$. Si el modelo fuera perfecto sería ruido
gris sin estructura; lo rojo y azul es lo que el modelo no describe, aquí sobre todo los brazos espirales, que no
están entre sus piezas.
\item \textbf{(d)} La imagen menos bulbo, disco y anillo: queda la barra sola (más los brazos), y se ve dónde
termina. La barra combinada (morado) encierra esa barra; la de Ferrers (verde) se extiende bastante más allá,
porque su $a_{\rm bar}$ es el radio donde el perfil llega a cero, no donde la barra deja de verse.
\end{itemize}
El ajuste se hace \emph{galaxia por galaxia}: es lento, necesita valores iniciales y elegir qué componentes
incluir, y en barras chicas falla con frecuencia. Por eso no forma parte del pipeline. Sirve para dos cosas: confirma
el ángulo de las isofotas ($133.8^\circ$ frente a $130.8^\circ$, y no los $120.5^\circ$ de Fourier) y, en el panel
(d), confirma que la barra combinada sigue siendo la que mejor delimita la barra visible.


\section{Resultados}\label{sec:resultados}

\subsection{Frente a Valeria: 448 barras}

Valeria midió 448 galaxias. 157 están en nuestra muestra DESI DR1 (las que se usaron para corregir el código);
las 291 restantes no tienen espectro DESI pero sí imagen Legacy, y se midieron con la posición y el $z$ de SDSS y
un recorte de tamaño $0.66\,R_{90}$ (el cociente mediano entre el radio de Legacy y el de SDSS en las 157). Son la
prueba justa: no se usaron para ajustar nada. La tabla~\ref{tab:valeria} y la figura~\ref{fig:valeria} resumen
la comparación (mediana del cociente nuestro/Valeria, semiancho 16--84\,\% y fracción a menos del 20\,\%).

\begin{table}[h]
\centering
\begin{tabular}{lccc}
\toprule
 & 157 (ajuste) & 291 (no usadas) & 448 \\
\midrule
Semilargo, isofotas & $1.01\pm0.10$ (89\,\%) & $1.01\pm0.11$ (88\,\%) & $1.01\pm0.10$ (88\,\%) \\
Semilargo, Fourier & $1.28\pm0.31$ (35\,\%) & $1.36\pm0.41$ (32\,\%) & $1.31\pm0.39$ (33\,\%) \\
Ancho, isofotas ($2b_{\rm iso}$) & $0.97\pm0.12$ & $0.97\pm0.12$ & $0.97\pm0.12$ \\
Ancho, combinada ($W_F$) & $0.47\pm0.10$ & $0.50\pm0.09$ & $0.49\pm0.10$ \\
Largo en kpc, isofotas & $1.05\pm0.10$ & $1.05\pm0.11$ & $1.05\pm0.11$ \\
\bottomrule
\end{tabular}
\caption{Cociente nuestro/Valeria (proyectado). Entre paréntesis, la fracción a menos del 20\,\%.}
\label{tab:valeria}
\end{table}

\begin{itemize}
\item \textbf{Largo}: coincide, y lo mismo en las 291 no usadas. Solo 7 de 448 quedan fuera de $0.6$--$1.5$.
\item \textbf{Ancho}: el de las isofotas coincide con el de Valeria porque los dos pasan por el centro e incluyen
el bulbo; el de Fourier es la mitad porque no lo incluye. Es otra definición, no un error.
\item \textbf{Largo en kpc}: $1.05$ en vez de $1.01$ porque nuestra escala (Planck18) es 3.4\,\% mayor que la de
Valeria; con la misma cosmología el cociente es el de los segundos de arco.
\item \textbf{Fourier} se pasa en largo ($1.3$, con el triple de dispersión): en barras chicas y discos inclinados
la fase se mantiene hacia los brazos.
\end{itemize}

\begin{figure}[t]
\centering
\includegraphics[width=\textwidth]{figures_barras/valeria.pdf}
\caption{Comparación con Valeria. (a) Semilargo por isofotas. (b) Ancho por isofotas (azul) y combinado
(morado). (c) Largo en kpc. Llenos: las 157 usadas para corregir el código; vacíos: las 291 no usadas.}
\label{fig:valeria}
\end{figure}

\subsection{DESI frente a SDSS}

La galaxia ejemplo medida con el mismo código en las dos imágenes (figura~\ref{fig:desi_sdss}): $\aeps = 15.4$ y
$15.0\as$, $\varepsilon_{\max} = 0.67$ en ambas, $a_\varphi = 11.8$ y $11.9\as$, $W_F = 5.6$ y $5.5\as$. El
método no depende de la imagen. La diferencia con Valeria en esta galaxia ($14.05\as$) es del método: con la misma
imagen SDSS, el largo por isofotas es 7\,\% mayor. Lo que sí depende de la imagen es la profundidad: las isofotas
llegan a $42.9\as$ en DESI y a $25.7\as$ en SDSS, lo que afecta la orientación del disco
(sección~\ref{sec:pendientes}).

\begin{figure}[t]
\centering
\includegraphics[width=\textwidth]{figures_barras/desi_sdss.pdf}
\caption{Galaxia ejemplo en (a) DESI y (b) SDSS, con las barras combinada (morado) y de Fourier (naranja).
(c) Elipticidad y (d) fase $\varphi_2$ en las dos imágenes.}
\label{fig:desi_sdss}
\end{figure}

Para separar el método de los datos en toda la muestra, se midieron las 448 galaxias de Valeria también en
las imágenes SDSS, con el mismo código (\code{measure\_bars.py --survey sdss}; ninguna falla). Cociente mediano,
mitad del rango 16--84\,\% y fracción a menos del 20\,\%:
\begin{center}
\begin{tabular}{lccc}
\toprule
 & Largo $a$ & Ancho Fourier $W$ & Ancho isofota $2b_{\rm iso}$ \\
\midrule
Método: SDSS nuestro / Valeria & $1.01\pm0.11$, 87\,\% & $0.48\pm0.10$ & $0.97\pm0.12$, 85\,\% \\
Datos: DESI nuestro / SDSS nuestro & $1.03\pm0.04$, 95\,\% & $1.01\pm0.08$, 94\,\% & $1.00\pm0.07$, 95\,\% \\
Total: DESI nuestro / Valeria & $1.01\pm0.10$, 88\,\% & $0.48\pm0.10$ & $0.97\pm0.12$, 86\,\% \\
\bottomrule
\end{tabular}
\end{center}
Casi toda la diferencia con Valeria es de método: con la misma imagen la dispersión es del 11\,\%, y entre
imágenes, del 4\,\%. DESI da barras un 3\,\% más largas que SDSS, por su mayor profundidad y mejor seeing. En la
mediana los dos métodos coinciden en el largo; el ancho de Fourier es la mitad del de Valeria en las dos
imágenes, porque no incluye el bulbo (figura~\ref{fig:metodo_datos}).

\begin{figure}[t]
\centering
\includegraphics[width=\textwidth]{figures_barras/metodo_datos.png}
\caption{Semilargo de las 448 barras. (a) Método: nuestro en SDSS frente a Valeria. (b) Datos: nuestro en DESI
frente a nuestro en SDSS. (c) Total: nuestro en DESI frente a Valeria. Línea: 1:1.}
\label{fig:metodo_datos}
\end{figure}

\subsection{La banda: $g$, $r$, $z$ y su suma}\label{sec:bandas}

Hasta aquí todo se midió en $r$. Las imágenes Legacy existen también en $g$ ($\sim$480\,nm, verde-azul) y $z$
($\sim$920\,nm, infrarrojo cercano), con el mismo campo, así que se midieron las 448 galaxias de Valeria con el
código sin cambios en $g$, $r$, $z$, en la suma de los tres FITS ($g+r+z$) y en el JPG de color pasado a gris.
En cada caso la barra se lee igual: el semilargo $a$ es el radio de la isofota donde la elipticidad es máxima
($\aeps$, sección~\ref{sec:isofotas}) y el ancho $W$ sale de Fourier (sección~\ref{sec:fourier}); lo único que
cambia es la imagen de entrada. La figura~\ref{fig:bandas} lo muestra en la galaxia ejemplo.

\begin{figure}[t]
\centering
\includegraphics[width=\textwidth]{figures_barras/bandas.pdf}
\caption{Galaxia ejemplo en (a) $g$, (b) $r$, (c) $z$ y (d) la suma $g+r+z$. La caja de cada color es la barra
medida en esa imagen: largo por isofotas ($\aeps$), ancho por Fourier. La franja vertical es $\pm10$\,\% del
largo, la dispersión frente a Valeria en las 448; el círculo a trazos, el semilargo de Valeria ($14.05\as$).
(e) Elipticidad de las isofotas en cada imagen; las líneas verticales marcan $\aeps$. (f) Resumen: largo con
$\pm10$\,\%; en gris, $\pm20$\,\% de Valeria.}
\label{fig:bandas}
\end{figure}

En el ejemplo, la barra se acorta del azul al infrarrojo: $16.2\as$ en $g$, $15.4\as$ en $r$, $14.0\as$ en $z$
y $14.7\as$ en la suma, es decir $1.15$, $1.10$, $0.99$ y $1.04$ veces el valor de Valeria. Las cuatro caen
dentro del 20\,\% y sus barras de error se solapan. Las cuatro curvas de elipticidad tienen la misma meseta
(de $13$ a $17.5\as$) y caen en el mismo radio; lo que cambia es en qué punto de la meseta queda el máximo, y el
nivel: $g$ es la más alargada.

\paragraph{Por qué la barra se acorta hacia el infrarrojo.} La barra es una estructura de estrellas viejas en
órbitas alargadas. El gas que la barra empuja hacia adentro forma estrellas sobre todo en el centro y en sus
extremos, donde nacen los brazos o hay un anillo; a lo largo de la barra la formación estelar suele estar
suprimida. Los brazos también contienen estrellas viejas (son ondas de densidad del disco), pero en ellos el gas
comprimido forma estrellas jóvenes, azules y brillantes. En $g$ esos extremos y el arranque de los brazos
resaltan, y además el polvo oscurece partes de la barra: las isofotas siguen alargadas y con el mismo ángulo más
allá del fin de la barra, y $\aeps$ crece. En $z$ domina la luz de las estrellas viejas, que sigue la masa: el
arranque de los brazos pesa menos y la elipticidad se corta antes. Por esta razón las barras se miden de
preferencia en el infrarrojo.

\paragraph{En las 448.} La tabla resume las 427 galaxias con las cinco medidas (en 21 falla la lectura del JPG):

\begin{center}
\begin{tabular}{lccccc}
\toprule
 & $r$ & $g$ & $z$ & suma & JPG \\
\midrule
A menos del 20\,\% & 88.5\,\% & 83.4\,\% & 89.0\,\% & 87.8\,\% & 38.6\,\% \\
A menos del 10\,\% & 65.6\,\% & 56.0\,\% & 69.1\,\% & 68.4\,\% & 25.3\,\% \\
Mediana / Valeria & 1.01 & 1.04 & 0.99 & 1.00 & 0.74 \\
Semiancho 16--84\,\% & 0.10 & 0.12 & 0.10 & 0.10 & 0.31 \\
Arregla / rompe frente a $r$ & -- & 15 / 37 & 18 / 16 & 7 / 10 & 16 / 229 \\
\bottomrule
\end{tabular}
\end{center}

\begin{itemize}
\item $r$, $z$ y la suma son equivalentes: la diferencia (un punto porcentual) es menor que el error de la
fracción con 427 galaxias ($\sim$1.5\,\%). $z$ tiene la mejor justificación física y la menor dispersión; $g$ se
pasa (mediana $1.04$) por lo explicado arriba.
\item El JPG es mucho peor, y no por las bandas: el estiramiento no lineal a 8 bits satura el centro, se pierden
las isofotas internas y la barra sale corta (mediana $0.74$). Las tres bandas ya están, sin saturar, en la suma de
los FITS.
\item De las 49 galaxias que fallan en $r$ (28 cortas, 21 se pasan), en 40 la suma da lo mismo que $r$ (a menos
del 10\,\%): la diferencia con Valeria no es de la banda ni del ruido, sino de dónde se pone el fin de la barra.
\item La dispersión entre bandas ($\sim$3\,\% en las que coinciden con Valeria, medida en un subconjunto de 102 en $griz$) es menor que la diferencia con
Valeria ($\sim$10\,\%): no sirve como incertidumbre. Se usa $\pm10$\,\%, la dispersión calibrada en las 448.
\end{itemize}

\paragraph{Solo $z$ frente a solo $r$.} Con el mismo método y en las 448, medir solo en $z$ da el mismo
acuerdo con Valeria que medir solo en $r$ (figura~\ref{fig:r_z}):
\begin{center}
\begin{tabular}{lcccc}
\toprule
 & A menos del 20\,\% & A menos del 10\,\% & Mediana / Valeria & Semiancho 16--84\,\% \\
\midrule
Solo $r$ & 88.4\,\% & 65.2\,\% & 1.01 & 0.10 \\
Solo $z$ & 88.6\,\% & 69.2\,\% & 0.99 & 0.10 \\
\bottomrule
\end{tabular}
\end{center}
Cambiar de $r$ a $z$ arregla 20 galaxias y rompe 19. Separadas por tamaño de la barra o por redshift
(figura~\ref{fig:r_z}b,c), las dos bandas difieren en 2--3 puntos, menos que el error de cada grupo
($\pm$3--5 puntos): $z$ no empeora en las galaxias lejanas, que era la duda por su menor profundidad, y $r$ gana
por poco en las barras grandes y cercanas, donde brazos y anillos se ven mejor. Elegir una u otra es una
convención: $z$ tiene la mejor justificación física y algo más de acuerdo fino; $r$ es la banda con la que se hizo
todo lo anterior y es más profunda.

Lo que sí aporta medir las dos es un aviso: cuando $|r/z - 1| > 0.15$ (42 galaxias), falla en $r$ el 45\,\%,
frente al 12\,\% de toda la muestra. Son galaxias con dos estructuras alargadas (un bulbo o lente alargado
dentro de la barra, o brazos y anillo fuera) en las que cada banda elige una. Se propone medir en $r$, medir
también en $z$ y marcar esas galaxias.

\begin{figure}[h!]
\centering
\includegraphics[width=\textwidth]{figures_barras/r_z.pdf}
\caption{Solo $r$ frente a solo $z$ en las 448. (a) Cociente con Valeria en $z$ contra el de $r$; líneas
punteadas: $\pm20$\,\%. (b) Fracción a menos del 20\,\% de Valeria por semilargo de Valeria y (c) por redshift,
con error binomial; morado: $r$, naranja: $z$.}
\label{fig:r_z}
\end{figure}

\subsection{Galaxias sin referencia}

La figura~\ref{fig:galeria} muestra la barra combinada en 5 de las 291 no usadas y en 5 galaxias ciegas de DESI
DR1, elegidas al azar. En las ciegas las barras son mucho más chicas ($2$--$4\as$): la muestra de Valeria está
sesgada a galaxias grandes y cercanas, y la mediana de nuestra muestra es $z \approx 0.06$. En esas barras la
medida es posible pero está cerca del límite de resolución.

\begin{figure}[t]
\centering
\includegraphics[width=\textwidth]{figures_barras/galeria.pdf}
\caption{Barra combinada (morado) con la barra sobre $x$. Arriba: 5 de las 291 de Valeria no usadas (círculo: su
semilargo). Abajo: 5 galaxias ciegas de DESI DR1, sin referencia.}
\label{fig:galeria}
\end{figure}


\section{Errores encontrados y cómo se corrigieron}\label{sec:errores}

La primera corrida sobre las 157 galaxias en común con Valeria dio mediana correcta pero muchos casos lejísimos.
No eran fallas del método sino del código (figura~\ref{fig:problemas}):
\begin{enumerate}
\item \textbf{Centro en una estrella} (13 galaxias). El centro era el píxel más brillante de toda la región de la
galaxia; si una estrella la tocaba, ganaba (a $33.9\as$ del núcleo en el panel a). Ahora se busca a menos de
10\,px del centro del recorte, que está en la posición DESI.
\item \textbf{Isofotas congeladas} ($\sim 22$ galaxias). \code{fit\_image} de \code{photutils} congela la geometría
de todos los radios mayores cuando dos isofotas seguidas tienen gradiente $\ge 0$, cosa común en barras y
anillos (panel b: llega a $3.9\as$ en vez de $26.3\as$). Ahora cada radio se ajusta por separado partiendo de la
última elipse buena.
\item \textbf{Isofotas pegadas al inicio} (16 galaxias). Cuando un radio no converge en 50 iteraciones (código 2),
\code{photutils} devuelve la geometría inicial como si fuera un ajuste. Con el inicio en $\varepsilon = 0.2$,
$\mathrm{PA} = 0$, eso crea un salto falso de $\varepsilon$ que parece el fin de una barra (panel c: 7 isofotas con
$\varepsilon = 0.2$ exacto entre $3.9$ y $5.3\as$). Ahora el ajuste parte de la forma de la luz, medida con sus
segundos momentos, y un código 2 cuya $\varepsilon$ no se movió se descarta.
\item \textbf{Pico equivocado}. Con isofotas que ya llegan al disco externo, el máximo global de $\varepsilon$ caía
a veces allí. Se usa el primer pico que cumple Marinova \& Jogee.
\end{enumerate}
La figura~\ref{fig:versiones} muestra el efecto de cada corrección sobre el largo en las 157: la fracción a menos
del 20\,\% de Valeria pasa de 69 a 89\,\% y las galaxias fuera de $0.6$–$1.5$ de 24 a 3. Promediar el PA sobre la
barra (v4) no cambia el largo; el ángulo se mueve $1.4^\circ$ en mediana.

\begin{figure}[t]
\centering
\includegraphics[width=\textwidth]{figures_barras/problemas.pdf}
\caption{Tres errores del primer código. (a) Centro viejo en una estrella (naranja) y nuevo (azul).
(b) Elipticidad con \code{fit\_image} (congelada a partir de $3.9\as$) y con el ajuste radio por radio. (c) Ajuste
desde $\varepsilon = 0.2$, PA $= 0$ (naranja), con las isofotas que nunca se movieron (cuadrados), y el actual.}
\label{fig:problemas}
\end{figure}

\begin{figure}[t]
\centering
\includegraphics[width=0.85\textwidth]{figures_barras/versiones.pdf}
\caption{Largo por isofotas en las 157 galaxias en común con Valeria, versión por versión. v0: primer código;
v1: centro cerca del recorte, ajuste radio por radio, primer pico; v2: isofotas pegadas (criterio parcial) y ancho
combinado; v3: inicio en los momentos de la luz y criterio completo; v4: PA promediado. (a) Mediana y semiancho
16--84\,\% de $\aeps/a_{\rm ref}$. (b) Fracción a menos del 20\,\% (azul) y número fuera de $0.6$--$1.5$ (naranja).}
\label{fig:versiones}
\end{figure}


\section{Limitaciones y pendientes}\label{sec:pendientes}

\begin{itemize}
% \item \textbf{Resolución.} Con seeing de $\sim 1.2\as$, una barra de $2$--$3\as$ está a pocos FWHM. Hay que fijar,
% con números propios, desde qué tamaño aparente la medida es confiable ($a_{\min} = 2\as$ es provisorio).
% \item \textbf{Deproyección.} La orientación del disco se toma de las isofotas externas ($a \ge 0.7\,a_{\text{último}}$),
% que en imágenes poco profundas o con anillos trazan el anillo y no el disco. Hasta resolverlo, todo es proyectado.
\item \textbf{Resolución.} Con seeing de $\sim 1.2\as$, una barra de $2$--$3\as$ está a pocos FWHM; $2\as$ sigue
siendo provisorio. El mínimo de 1\,kpc (sección~\ref{sec:metodo}) solo resuelve el problema contrario, el de las
galaxias cercanas.
% \item \textbf{Errores.} Falta el error de cada medida: para $a$, el tramo donde $\varepsilon$ es compatible con el
% pico; para $W$, la dispersión entre cortes; la deproyección, por Monte Carlo.
\item \textbf{Errores en el código.} El error de la meseta (sección~\ref{sec:error}) está probado pero no en
\code{src/}; falta decidir cómo se lleva a $W$.
\item \textbf{Salida.} Falta el catálogo corto (TARGETID, RA, DEC, $z$, $a$, $L$, $W$ con errores, PA y una
columna de calidad) y el visualizador por galaxia.
\item \textbf{Lista sin \code{SHAPE\_R}.} Elegir el lado del recorte sin Tractor (sección~\ref{sec:lista}).
\item \textbf{Máscaras.} Las estrellas superpuestas a la galaxia y los artefactos (trazas de satélite) no se
enmascaran; la mediana los atenúa. Se probará GLIDErS (Erostegui), que enmascara estrellas superpuestas y
artefactos.
\item \textbf{Ángulo.} En 14 de 157 galaxias isofotas y Fourier difieren más de $10^\circ$; a ojo y según la
descomposición, el que falla es Fourier. Se usará como bandera de revisión.
\item \textbf{Definición del fin de la barra.} El largo por isofotas coincide con Valeria, pero el largo es una
definición (no hay borde nítido, sección~\ref{sec:contexto}): en la galaxia ejemplo los dos dan $14$--$15\as$ y a
ojo se ve $\sim 11\as$, que es lo que da la fase de Fourier. La cota superior $a_{\rm up}$ se calculó en las
versiones de validación (\code{checks/}) pero no está en el código simplificado.
\item \textbf{Tamaño normalizado y fuerza.} Para comparar pares y control hace falta el largo normalizado por el
tamaño de la galaxia y comparar a igual masa (sección~\ref{sec:contexto}); falta elegir el radio de normalización.
\end{itemize}


\section{Del documento al código}\label{sec:codigo}

\begin{table}[h]
\centering
\small
\begin{tabularx}{\textwidth}{lX}
\toprule
Paso & Dónde \\
\midrule
Muestra (GZ DESI $p_{\rm bar} > 0.5$, DR1 a $< 1\as$, $z < 0.1$) & \code{src/sample.py}, \code{scripts/build\_sample.py} \\
Recortes FITS y JPG, DESI y SDSS & \code{src/cutout.py} \\
Cielo, máscara, centro, isofotas radio por radio & \code{src/isophotes.py} \\
Barra en $\varepsilon(a)$, PA$(a)$; kpc \eqref{eq:kpc} & \code{src/bar.py} \\
Modos de Fourier \eqref{eq:fourier} y ancho \eqref{eq:wf} & \code{src/fourier.py} \\
Todo junto: \code{measure\_fits}, \code{measure\_image}, \code{measure\_galaxy}; barra combinada \eqref{eq:comb} & \code{src/measure.py} \\
Una muestra entera, también por bloques en un cluster & \code{scripts/measure\_bars.py}, \code{scripts/merge\_bars.py} \\
Comparación con Valeria & \code{scripts/build\_reference\_sample.py}, \code{tools/compare\_reference.py} \\
Galerías; el método paso a paso & \code{tools/plot\_gallery.py}; \code{notebooks/Barra\_ejemplo.ipynb} (\code{tools/steps.py}) \\
Validación: deproyección \eqref{eq:deproj}, largo por la fase, descomposición 2D \eqref{eq:chi2}--\eqref{eq:ferrers} & \code{checks/} \\
\bottomrule
\end{tabularx}
\end{table}

\clearpage

\section*{Referencias}
\small
Las de la sección~\ref{sec:contexto} vienen de las diapositivas de Cuomo; verificar en ADS antes de usarlas en un paper.\\[2pt]
Aguerri, J. A. L., Muñoz-Tuñón, C., Varela, A. M., \& Prieto, M. 2000, A\&A, 361, 841.\\
Aguerri, J. A. L., Méndez-Abreu, J., \& Corsini, E. M. 2009, A\&A, 495, 491.\\
Alonso, S., Vera-Rueda, M., Coldwell, G., Duplancic, F., \& Mesa, V. 2024, arXiv:2409.00194.\\
Cuomo, V. 2025, Galactic bars and their influence on galaxy evolution, clases 2--5, Formación y Evolución de
Galaxias, Universidad de La Serena.\\
Cuomo, V., et al. 2020, A\&A, 641, A111.\\
Elmegreen, B. G., \& Elmegreen, D. M. 1985, ApJ, 288, 438 (verificar en ADS).\\
de Vaucouleurs, G. 1959, Handbuch der Physik, 53, 275.\\
Erwin, P. 2005, MNRAS, 364, 283.\\
Erwin, P. 2015, ApJ, 799, 226.\\
Gadotti, D. A. 2008, MNRAS, 384, 420.\\
Herrera-Endoqui, M., Díaz-García, S., Laurikainen, E., \& Salo, H. 2015, A\&A, 582, A86.\\
Hilmi, T., et al. 2020, MNRAS, 497, 933.\\
Hubble, E. P. 1926, ApJ, 64, 321.\\
Jedrzejewski, R. I. 1987, MNRAS, 226, 747.\\
Marinova, I., \& Jogee, S. 2007, ApJ, 659, 1176.\\
Ohta, K., Hamabe, M., \& Wakamatsu, K. 1990, ApJ, 357, 71 (verificar en ADS).\\
Martínez-Valpuesta, I., et al. 2017, MNRAS, 464, 1502.\\
Peng, C. Y., Ho, L. C., Impey, C. D., \& Rix, H.-W. 2002, AJ, 124, 266.\\
Sheth, K., et al. 2010, PASP, 122, 1397.\\
Walmsley, M., et al. 2023, MNRAS, 526, 4768.\\
Wozniak, H., \& Pierce, M. J. 1991, A\&AS, 88, 325.

\end{document}

